% from ega3-6-fr.tex
\oldpage[III]{167}
\[
\begin{split}
  {}^{(e)}\mathscr E_{pq}^2
  =\bigoplus_{q'+q''=q}
  \mathscr{T}\!or_p^S\Bigl(
  &\mathfrak{Tor}_{q'}^S
  (f_1,\ldots,f_m;
   \mathscr P_\bullet^{(1)},\ldots,\mathscr P_\bullet^{(m)}),\\
  &\mathscr H_{q''}(\mathscr Q'_\bullet)
  \Bigr),
\end{split}
\]
\[
\begin{split}
  {}^{(f)}\mathscr E_{pq}^2
  =\bigoplus_{q_1+q_2+\cdots+q_{m+1}=q}
  \mathscr{T}\!or_p^{S'}\Bigl(
  &\mathfrak{Tor}_{q_1}^S
  (f_1,1_{S'};\mathscr P_\bullet^{(1)},\mathscr O_{S'}),\ldots,\\
  &\mathfrak{Tor}_{q_m}^S
  (f_m,1_{S'};\mathscr P_\bullet^{(m)},\mathscr O_{S'}),
  \mathscr H_{q_{m+1}}(\mathscr Q'_\bullet)
  \Bigr),
\end{split}
\]
\[
\begin{split}
  {}^{(f')}\mathscr E_{pq}^2
  =\bigoplus_{q_1+\cdots+q_m=q}
  \mathfrak{Tor}_p^{S'}\Bigl(
  &f'_1,\ldots,f'_m,1_{S'};\\
  &\mathscr{T}\!or_{q_1}^S
  (\mathscr P_\bullet^{(1)},\mathscr O_{S'}),\ldots,
  \mathscr{T}\!or_{q_m}^S
  (\mathscr P_\bullet^{(m)},\mathscr O_{S'}),
  \mathscr Q'_\bullet
  \Bigr).
\end{split}
\]
\end{proposition}

\begin{proof}
La suite $(e)$ n'est autre que la suite d'associativité de (6.8.2) pour
$r=m$. Pour définir les deux autres suites spectrales, on va encore se limiter
au cas où $S$, $S'$ et les $Y^{(i)}$ sont affines, le passage au cas général
se faisant par les méthodes de (6.7.1) et (6.7.2) et étant laissé au lecteur.
Nous démontrerons donc le
\end{proof}

\begin{corollary}[6.9.4]
Soient $A$ un anneau, $A'$ une $A$-algèbre,
$S=\operatorname{Spec}(A)$, $S'=\operatorname{Spec}(A')$,
$X^{(i)}$ ($1\leq i\leq m$) des $S$-schémas quasi-compacts et, pour chaque
$i$, soit $X'^{(i)}=X_{(S')}^{(i)}$, qui est un $S'$-schéma quasi-compact.
Pour chaque $i$, soit $\mathscr P_\bullet^{(i)}$ un complexe de
$\mathscr O_{X^{(i)}}$-Modules quasi-cohérents; soit enfin $Q'_\bullet$ un
complexe de $A'$-modules, ces complexes étant limités inférieurement. Il
existe trois foncteurs spectraux biréguliers en les
$\mathscr P_\bullet^{(i)}$ et en $Q'_\bullet$, ayant pour aboutissement commun
\[
  \mathbf{Tor}_\bullet^S
  (X^{(1)},\ldots,X^{(m)},S';
   \mathscr P_\bullet^{(1)},\ldots,\mathscr P_\bullet^{(m)},
   \widetilde{Q'_\bullet})
\]
et dont les termes $E_2$ sont respectivement
\[
\begin{split}
  {}^{(e)}E_{pq}^2
  =\bigoplus_{q'+q''=q}
  \operatorname{Tor}_p^A\Bigl(
  &\mathbf{Tor}_{q'}^S
  (X^{(1)},\ldots,X^{(m)};
   \mathscr P_\bullet^{(1)},\ldots,\mathscr P_\bullet^{(m)}),\\
  &H_{q''}(Q'_\bullet)
  \Bigr),
\end{split}
\]
\[
\begin{split}
  {}^{(f)}E_{pq}^2
  =\bigoplus_{q_1+q_2+\cdots+q_{m+1}=q}
  \operatorname{Tor}_p^{A'}\Bigl(
  &\mathbf{Tor}_{q_1}^S
  (X^{(1)},S';\mathscr P_\bullet^{(1)},\mathscr O_{S'}),\ldots,\\
  &\mathbf{Tor}_{q_m}^S
  (X^{(m)},S';\mathscr P_\bullet^{(m)},\mathscr O_{S'}),
  H_{q_{m+1}}(Q'_\bullet)
  \Bigr),
\end{split}
\]
\[
\begin{split}
  {}^{(f')}E_{pq}^2
  =\bigoplus_{q_1+\cdots+q_m=q}
  \mathbf{Tor}_p^{S'}\Bigl(
  &X'^{(1)},\ldots,X'^{(m)},S';\\
  &\mathscr{T}\!or_{q_1}^S
  (\mathscr P_\bullet^{(1)},\mathscr O_{S'}),\ldots,
  \mathscr{T}\!or_{q_m}^S
  (\mathscr P_\bullet^{(m)},\mathscr O_{S'}),
  \widetilde{Q'_\bullet}
  \Bigr).
\end{split}
\]
\end{corollary}

\begin{proof}
Nous ne reviendrons pas sur le premier de ces foncteurs spectraux, qui a été
traité dans (6.8.3) et n'est inclus ici que pour mémoire. Pour définir les
autres, considérons pour chaque $i$ un recouvrement ouvert affine fini
$\mathfrak U^{(i)}$ de $X^{(i)}$ et, si
$u_i:X'^{(i)}\to X^{(i)}$ est la projection canonique, le recouvrement ouvert
affine fini correspondant
$\mathfrak U'^{(i)}=u_i^{-1}(\mathfrak U^{(i)})$.

En vertu de (6.6.6),
$\mathbf{Tor}_\bullet^S(X^{(1)},\ldots,X^{(m)},S';
\mathscr P_\bullet^{(1)},\ldots,\mathscr P_\bullet^{(m)},
\widetilde{Q'_\bullet})$ s'obtient en prenant pour $1\leq i\leq m$ une
résolution projective de Cartan-Eilenberg
$M_{\bullet\bullet\bullet}^{(i)}$ de
$L_{\bullet\bullet}^{(i)}=C^\bullet(\mathfrak U^{(i)},
\mathscr P_\bullet^{(i)})$ (au sens de (0, 11.7.1)), considérant le
tricomplexe
\[
  M_{\bullet\bullet\bullet}
  =M_{\bullet\bullet\bullet}^{(1)}\otimes_A
   M_{\bullet\bullet\bullet}^{(2)}\otimes_A\cdots\otimes_A
   M_{\bullet\bullet\bullet}^{(m)}\otimes_A Q'_\bullet
\]
(où $Q'_\bullet$ est considéré comme un tricomplexe dont les deux derniers
degrés se réduisent à $0$), et en en prenant l'homologie. Si l'on pose
$M_{\bullet\bullet\bullet}'^{(i)}=
M_{\bullet\bullet\bullet}^{(i)}\otimes_A A'$, on a (en se rappelant que
$Q'_\bullet$ est un complexe de $A'$-modules)
\[
  M_{\bullet\bullet\bullet}
  =M_{\bullet\bullet\bullet}'^{(1)}\otimes_{A'}
   M_{\bullet\bullet\bullet}'^{(2)}\otimes_{A'}\cdots\otimes_{A'}
   M_{\bullet\bullet\bullet}'^{(m)}\otimes_{A'} Q'_\bullet.
\]
Or, considérons chacun des complexes
$M_{\bullet\bullet\bullet}'^{(i)}$ comme un complexe simple (pour son degré
total) et notons que ce complexe est formé de $A'$-modules projectifs; il
résulte de (6.3.7) (étendu à un nombre quelconque de complexes) que
$H_\bullet(M_{\bullet\bullet\bullet})$ est aussi égal à
\[
  \operatorname{Tor}_\bullet^{A'}
  (M_{\bullet\bullet\bullet}'^{(1)},\ldots,
   M_{\bullet\bullet\bullet}'^{(m)},Q'_\bullet);
\]
c'est donc (6.5.15) l'aboutissement d'une suite spectrale ayant les propriétés
de régularité voulues (les trois degrés de
$M_{\bullet\bullet\bullet}^{(i)}$ étant limités inférieurement lorsque
$\mathscr P_\bullet^{(i)}$ est limité inférieurement) et dont le terme $E_2$
est donné par
\[
  E_{pq}^2
  =\bigoplus_{q_1+\cdots+q_{m+1}=q}
  \operatorname{Tor}_p^{A'}
  (H_{q_1}(M_{\bullet\bullet\bullet}'^{(1)}),\ldots,
   H_{q_m}(M_{\bullet\bullet\bullet}'^{(m)}),H_{q_{m+1}}(Q'_\bullet)).
\]

